\documentclass[11pt]{article}
\usepackage[T1]{fontenc}
\usepackage[utf8]{inputenc}
\usepackage[a4paper,margin=1in]{geometry}
\usepackage{amsmath}
\usepackage{amssymb}
\usepackage{booktabs}
\usepackage{hyperref}
\usepackage{xcolor}

\hypersetup{
  colorlinks=true,
  linkcolor=blue,
  urlcolor=blue,
  pdftitle={RefPlay Architecture Notes},
  pdfauthor={OpenAI Codex}
}

\title{RefPlay: Architecture and Training Objective}
\author{}
\date{\today}

\begin{document}
\maketitle

\section{Overview}

RefPlay is an online preference-learning pipeline for math problems.
For each prompt $x$, two models produce answers:
\begin{itemize}
  \item the \textbf{actor} $\pi_{\theta}$, which is trainable;
  \item the \textbf{reference} $\pi_{\mathrm{ref}}$, which is frozen.
\end{itemize}

The actor and the reference answer the same prompt, both answers are scored,
and the actor is updated from the resulting chosen/rejected pair.

\section{Architecture}

The pipeline has three main stages.

\subsection*{1. Generation}

Given a prompt $x$:
\begin{align*}
y_{\mathrm{actor}} &\sim \pi_{\theta}(\cdot \mid x), \\
y_{\mathrm{ref}} &\sim \pi_{\mathrm{ref}}(\cdot \mid x).
\end{align*}

The current implementation uses:
\begin{itemize}
  \item \textbf{actor initialization:} \texttt{Qwen/Qwen2.5-Math-1.5B-Instruct},
  \item \textbf{reference initialization:} \texttt{Qwen/Qwen2.5-Math-1.5B-Instruct}.
\end{itemize}

The actor is optimized during training; the reference is never updated.

\subsection*{2. Rewarding}

Each response is scored independently:
\begin{align*}
s_{\mathrm{actor}} &= r(x, y_{\mathrm{actor}}), \\
s_{\mathrm{ref}} &= r(x, y_{\mathrm{ref}}).
\end{align*}

The current reward is a dense GSM8K-oriented reward:
\begin{itemize}
  \item full reward for a mathematically correct final answer when exact verification is available;
  \item bonus for producing a parseable final answer;
  \item bonus for using the expected \texttt{\#\#\#\#} final-answer format;
  \item small bonus for concise outputs;
  \item penalty for strong repetition.
\end{itemize}

This reward is denser than a pure binary correct/incorrect score, which makes
early training less sparse.

\subsection*{3. Preference Construction and Update}

The higher-scoring response becomes \texttt{chosen} and the other becomes
\texttt{rejected}. This creates an online preference pair:
\begin{align*}
y^{+} &= \arg\max \{ s_{\mathrm{actor}}, s_{\mathrm{ref}} \}, \\
y^{-} &= \arg\min \{ s_{\mathrm{actor}}, s_{\mathrm{ref}} \}.
\end{align*}

Then the actor computes policy log-probabilities for both responses, while the
frozen reference computes reference log-probabilities for both responses.

\section{Loss}

RefPlay uses a DPO-style objective. Let
\begin{align*}
\log \pi_{\theta}(y^{+} \mid x), \quad
\log \pi_{\theta}(y^{-} \mid x), \\
\log \pi_{\mathrm{ref}}(y^{+} \mid x), \quad
\log \pi_{\mathrm{ref}}(y^{-} \mid x)
\end{align*}
be sequence-level log-probabilities on the chosen and rejected responses.

The key logit is
\[
\Delta =
\bigl(\log \pi_{\theta}(y^{+}\mid x) - \log \pi_{\theta}(y^{-}\mid x)\bigr)
-
\bigl(\log \pi_{\mathrm{ref}}(y^{+}\mid x) - \log \pi_{\mathrm{ref}}(y^{-}\mid x)\bigr).
\]

The training loss is a sigmoid-style DPO loss:
\[
\mathcal{L}_{\mathrm{DPO}}
=
- \log \sigma(\beta \Delta),
\]
where $\beta > 0$ controls how strongly the actor is pushed to prefer the
chosen answer over the rejected one relative to the frozen reference.

\section{Interpretation}

This means the actor is trained to satisfy two goals at once:
\begin{itemize}
  \item beat the frozen reference on the task reward;
  \item increase the relative likelihood of the better answer over the worse one.
\end{itemize}

So the pipeline is best understood as:
\begin{quote}
online pairwise preference optimization with a frozen reference opponent.
\end{quote}

\section{Current Data Setting}

The current training and evaluation setup uses GSM8K-style math data:
\begin{itemize}
  \item data source: \texttt{openai/gsm8k},
  \item task type: math word problems,
  \item supervision target: final numeric answer,
  \item expected response pattern: reasoning followed by \texttt{\#\#\#\# answer}.
\end{itemize}

\section{What Is Frozen and What Is Updated}

\begin{center}
\begin{tabular}{lll}
\toprule
Component & Role & Updated? \\
\midrule
Actor & Generates answers and is optimized & Yes \\
Reference & Generates comparison answers and reference log-probs & No \\
Reward & Scores both answers & No gradients to model policy \\
\bottomrule
\end{tabular}
\end{center}

\section{Summary}

In short, RefPlay works as follows:
\begin{enumerate}
  \item sample a prompt;
  \item generate one answer from the actor and one from the frozen reference;
  \item score both answers with the reward function;
  \item form a chosen/rejected pair;
  \item update only the actor with a DPO-style loss against the frozen reference.
\end{enumerate}

\end{document}
